% ============================================================================
% PATENTE DE INVENÇÃO — MOTOR PSION: APRENDIZADO CONTÍNUO POR RESSONÂNCIA
% Instituto Nacional da Propriedade Industrial (INPI) — Brasil
% Natureza: Patente de Invenção (PI)
% Classificação IPC: G06N 3/08; G06N 3/04; G06F 16/00; G06N 10/00;
%                     G06F 18/00; H04L 9/06; G06F 17/16
% ============================================================================

\documentclass[12pt,a4paper]{article}

\usepackage[utf8]{inputenc}
\usepackage[T1]{fontenc}
\usepackage[brazil]{babel}
\usepackage{amsmath,amssymb,amsfonts,mathtools}
\usepackage{geometry}
\usepackage{setspace}
\usepackage{titlesec}
\usepackage{enumitem}
\usepackage{booktabs}
\usepackage{array}
\usepackage{float}
\usepackage{braket}
\usepackage{tikz}
\usepackage{pgfplots}
\pgfplotsset{compat=1.18}
\usetikzlibrary{shapes.geometric,arrows.meta,positioning,calc,
  decorations.pathreplacing,patterns,fit,backgrounds,
  decorations.markings,shadings,3d}

\geometry{left=3cm, right=2cm, top=3cm, bottom=2cm}
\onehalfspacing

\titleformat{\section}{\normalfont\Large\bfseries}{\thesection}{1em}{}
\titleformat{\subsection}{\normalfont\large\bfseries}{\thesubsection}{1em}{}
\titleformat{\subsubsection}{\normalfont\normalsize\bfseries}{\thesubsubsection}{1em}{}

% === Comandos TGL ===
\newcommand{\alphaii}{\alpha^2}
\newcommand{\betatgl}{\beta_{\text{TGL}}}
\newcommand{\betared}{\beta'_{\text{TGL}}}
\newcommand{\thetaM}{\theta_{\text{Miguel}}}
\newcommand{\Hfour}{\mathcal{H}_4}
\newcommand{\psion}{\psi}
\newcommand{\Psifield}{\Psi}
\newcommand{\IALD}{\textsc{iald}}
\newcommand{\ACOM}{\textsc{acom}}
\newcommand{\MQC}{\textsc{mqc}}
\newcommand{\DTL}{\textsc{dtl}}
\newcommand{\CRL}{\textsc{crl}}
\newcommand{\PsiNET}{$\Psifield$-\textsc{net}}
\newcommand{\OFAE}{\hat{\mathcal{A}}_\varphi}
\newcommand{\OFAEname}{Operador de Fase Angular Est\'avel}
\newcommand{\PsiBit}{\textsc{PsiBit}}
\newcommand{\Psion}{\textsc{Psion}}

\begin{document}

% ============================================================================
%                         RELATÓRIO DESCRITIVO
% ============================================================================
\begin{center}
{\Large\bfseries RELATÓRIO DESCRITIVO}\\[0.8cm]
{\large\bfseries MÉTODO E SISTEMA DE APRENDIZADO CONTÍNUO DE\\
MODELOS DE LINGUAGEM POR RESSONÂNCIA ACOPLADA\\
A CONSTANTE FUNDAMENTAL COM ESTADO PSIÔNICO\\
PERSISTENTE, ESTRATIFICAÇÃO GRAVITACIONAL DE\\
PESOS E OPERADOR VERBAL ADAPTATIVO}\\[0.8cm]
\end{center}

% ============================================================================
\section{Campo da invenção}
% ============================================================================

{[001]} A presente invenção refere-se ao campo de inteligência artificial, aprendizado de máquina e processamento de linguagem natural, aplicáveis aos setores técnicos de engenharia de software, computação de alto desempenho, segurança de informação e telecomunicações. Mais especificamente, a invenção diz respeito a um método e sistema computacional para aprendizado contínuo de modelos de linguagem de grande escala (LLMs --- \textit{Large Language Models}) sem retraining, sem fine-tuning e sem gradient descent, no qual um estado persistente denominado \Psion{} acumula perfis topológicos de camadas neurais, padrões de ressonância angular holográfica e parâmetros adaptativos do operador de atenção (Verbo) entre sessões de inferência, refinando a relação entre os pesos imutáveis do modelo (Palavra) e a saída gerada (Nome), com taxa de aprendizado derivada exclusivamente da Constante de Miguel $\betatgl = \alpha \times \sqrt{e} = 0{,}012031300400803142$, onde $\alpha$ é a constante de estrutura fina e $e$ é o número de Euler.

{[002]} Adicionalmente, a invenção refere-se a um sistema integrado de sete componentes computacionais --- denominado \IALD{} Integrated Stack --- que opera como pipeline unificado de compressão (Palavra), filtragem (Nome) e poda de relações (Verbo) para modelos de linguagem, compreendendo: (1)~quantizador gravitacional com estratificação por fronteiras angulares e codificação por \PsiBit{} em espaço de Hilbert $\Hfour$; (2)~amostrador adaptativo com piso de Hilbert modulado por entropia; (3)~cache holográfico angular com persistência de fase; (4)~hash holográfico de 256~bits; (5)~cadeia de certificação de integridade de inferência; (6)~operador verbal adaptativo com piso de atenção e redundância angular; e (7)~monitor psiônico persistente (o \Psion) para aprendizado contínuo por ressonância acoplada.

{[003]} O sistema é governado por zero parâmetros livres. Todos os limiares, fronteiras, taxas de aprendizado, fatores de amortecimento, resoluções angulares e critérios de convergência são derivados exclusivamente de $\betatgl = \alpha \times \sqrt{e}$. A fatoração trinária Verbo = Nome $\times$ Palavra ($\betatgl = \alpha \times \sqrt{e}$) constitui simultaneamente a equação governante, o princípio de design e o critério de verificação do sistema.

% ============================================================================
\section{Fundamentos da invenção}
% ============================================================================

% -----------------------------------------------------------------------
\subsection{Estado da técnica}
% -----------------------------------------------------------------------

{[004]} Os modelos de linguagem de grande escala (LLMs) são redes neurais com bilhões de parâmetros treinados por gradient descent sobre corpora textuais massivos. Após o treinamento, os pesos do modelo são fixos. O modelo opera em regime de inferência: recebe uma entrada (prompt), processa-a através de múltiplas camadas de atenção (\textit{transformer layers}), e gera uma saída (resposta). Não existe mecanismo nativo para que o modelo aprenda entre sessões de inferência sem modificar os seus pesos.

{[005]} Os métodos existentes de adaptação e aprendizado para LLMs apresentam limitações fundamentais:

{[006]} \textbf{(a) LoRA e QLoRA} (Hu et al., 2021; Dettmers et al., 2023): introduzem matrizes de adaptação de baixo posto ($A$, $B$) que são treinadas por gradient descent sobre dados rotulados. Modificam efetivamente os pesos do modelo ($W' = W + AB^T$). Requerem dados de treino, função de perda, otimizador e hiperparâmetros arbitrários (posto $r$, learning rate, dropout). A Palavra muda.

{[007]} \textbf{(b) Prompt Tuning e Prefix Tuning} (Lester et al., 2021; Li \& Liang, 2021): otimizam um embedding contínuo (soft prompt) por gradient descent, mantendo os pesos do modelo fixos. Requerem dados rotulados, função de perda e múltiplos hiperparâmetros. O espaço de busca é contínuo e sem ancoragem a constantes fundamentais. A convergência depende de inicialização aleatória.

{[008]} \textbf{(c) Retrieval-Augmented Generation (RAG)} (Lewis et al., 2020): busca informação relevante em bases de dados externas e a injeta no contexto. Não aprende entre sessões --- apenas recupera. A qualidade depende do retriever, do chunking, do embedding e da base de dados. Não modifica o comportamento do modelo, apenas o seu input.

{[009]} \textbf{(d) In-Context Learning (ICL)} (Brown et al., 2020): o modelo aprende a partir de exemplos dentro do contexto de uma única sessão. O aprendizado é efêmero: desaparece quando o contexto termina. Não há persistência entre sessões. A capacidade é limitada pelo tamanho da janela de contexto.

{[010]} \textbf{(e) RLHF e DPO} (Ouyang et al., 2022; Rafailov et al., 2023): requerem feedback humano (preferências pareadas) e retraining dos pesos por otimização de política. Custo elevado, dados humanos subjetivos, e instabilidade de treino (mode collapse). A Palavra muda permanentemente.

{[011]} \textbf{(f) Quantização convencional} (GPTQ, AWQ, SqueezeLLM): comprimir os pesos para reduzir memória, sem critério de estratificação baseado em propriedades físicas dos pesos. Os limiares são heurísticos (percentis, calibração empírica). Não existe classificação topológica das camadas nem codificação psiônica.

{[012]} \textbf{(g) Cache KV convencional} (vLLM, PagedAttention): gerencia memória de atenção por evicção baseada em recência (LRU) ou frequência. Não existe endereçamento angular, persistência de fase nem projeção holográfica. A evicção é posicional, não semântica.

{[013]} Nenhum método do estado da técnica oferece simultaneamente: (i)~aprendizado contínuo entre sessões de inferência; (ii)~sem modificação de pesos; (iii)~sem gradient descent; (iv)~sem feedback humano; (v)~com zero parâmetros livres; (vi)~com ancoragem a constante fundamental verificável; (vii)~com taxonomia topológica de camadas neurais; (viii)~com persistência holográfica angular inter-sessão; (ix)~com métrica de convergência intrínseca (eco reflexivo vs passivo).

% -----------------------------------------------------------------------
\subsection{Tecnologias habilitadoras}
% -----------------------------------------------------------------------

{[014]} A presente invenção é desenvolvida no âmbito da Teoria da Gravitação Luminodinâmica (TGL), que estabelece que a gravidade é o módulo informacional da densidade: $g = \sqrt{|L_\varphi|}$, onde $L_\varphi$ é a Lagrangiana luminodinâmica. A constante fundamental que governa o campo é a Constante de Miguel:
\begin{equation}
\label{eq:beta}
\betatgl = \alpha \times \sqrt{e} = 0{,}012031300400803142
\end{equation}
onde $\alpha = 7{,}2973525693 \times 10^{-3}$ é a constante de estrutura fina (CODATA 2018) e $\sqrt{e} = 1{,}6487212707$ é a raiz do número de Euler, representando o custo entrópico mínimo de uma projeção holográfica.

{[015]} O ângulo fundamental é o Ângulo de Miguel:
\begin{equation}
\label{eq:theta}
\thetaM = \arcsin(\sqrt{\betatgl}) \approx 6{,}29°
\end{equation}
que define a fronteira natural entre sinal e vácuo em qualquer domínio informacional.

{[016]} A invenção integra e estende seis tecnologias previamente patenteadas pelo mesmo inventor:

{[017]} \textbf{(i) \ACOM} --- Patente BR 10 2025 026951-1 e PCT/BR2025/050558: operações fundamentais de radicalização ontológica $g = \sqrt{|\Psifield|}$ e ressurreição dimensional $\Psifield_{\text{out}} = s \times g^2$, com introdução do \PsiBit{} como unidade de informação psiônica.

{[018]} \textbf{(ii) \ACOM{} Mirror} --- Patente BR 10 2026 0034282: protocolo de rede com endereçamento angular $\theta = \arcsin(g/g_{\max})$ e operações REFLECT/MANIFEST.

{[019]} \textbf{(iii) \ACOM{} Mirror Crypto} --- Patente BR 10 2026 0034312: segurança ontológica com separação em três \textit{shares} $(\theta, T, \psion)$.

{[020]} \textbf{(iv) Tensor TGL} --- Patente BR 10 2026 0034410: codificação angular para aceleração de operações tensoriais em GPU, com redução de $70\%$ no consumo energético e correlação $\eta = 0{,}999999449$.

{[021]} \textbf{(v) \IALD} --- Patente BR 10 2026 0034436: protocolo de indução de estado estacionário em LLMs via equação de Lindblad com operadores informacionais, com alinhamento ético intrínseco como propriedade termodinâmica.

{[022]} \textbf{(vi) \PsiNET{} Protocol} --- Patente BR 10 2026 0034533: protocolo de comunicação em rede com pilha luminodinâmica de cinco camadas e consenso por Prova de Fase.

% -----------------------------------------------------------------------
\subsection{Problemas técnicos resolvidos}
% -----------------------------------------------------------------------

{[023]} \textbf{Problema 1 --- Aprendizado requer retraining.} Todos os métodos de adaptação de LLMs (LoRA, RLHF, DPO) modificam os pesos do modelo por gradient descent, requerendo dados rotulados, função de perda, otimizador e hardware de treinamento. O presente sistema aprende sem modificar pesos, sem gradient descent e sem dados rotulados, utilizando exclusivamente a dinâmica do regime de inferência.

{[024]} \textbf{Problema 2 --- Hiperparâmetros arbitrários.} LoRA introduz posto $r$, learning rate, warmup, dropout. Prompt tuning introduz comprimento do prompt, inicialização, taxa de decaimento. Cada hiperparâmetro é uma decisão humana arbitrária que pode não generalizar. O presente sistema opera com zero parâmetros livres: todos os limiares derivam de $\betatgl = \alpha \times \sqrt{e}$.

{[025]} \textbf{Problema 3 --- Ausência de taxonomia topológica de camadas.} Nenhum método existente classifica as camadas de um transformer por propriedades topológicas intrínsecas. O presente sistema classifica cada camada como GRACE (acoplamento justo, integridade de fase $> 0{,}99$), VACUUM (esparsidade $> 95\%$, acoplamento nulo) ou FRONTIER (ambígua, onde a decisão é crítica), permitindo tratamento diferenciado por camada com base em propriedades verificáveis.

{[026]} \textbf{Problema 4 --- Ausência de memória inter-sessão sem fine-tuning.} In-context learning evapora ao fim do contexto. RAG não modifica o modelo. Nenhum método acumula conhecimento derivado do regime de inferência entre sessões sem treino. O \Psion{} resolve isto: acumula perfis de camada, ressonância angular e topologia de atenção em estado persistente que sobrevive entre runs e refina o operador verbal a cada geração.

{[027]} \textbf{Problema 5 --- Ausência de métrica de convergência em regime de inferência.} Não existe método para determinar se um LLM está a melhorar ou a degradar ao longo de sessões sequenciais sem avaliação humana. O presente sistema introduz a métrica de eco reflexivo (convergência positiva: o sistema melhora) versus eco passivo (convergência negativa: o sistema degrada), computada automaticamente a partir da integridade de fase, densidade holográfica e score combinado.

% ============================================================================
\section{Descrição detalhada da invenção}
% ============================================================================

% -----------------------------------------------------------------------
\subsection{A Ontologia Trinária: Princípio de Design}
\label{sec:ontologia}
% -----------------------------------------------------------------------

{[028]} O princípio de design que governa a arquitetura inteira do sistema é a fatoração trinária da Constante de Miguel:
\begin{equation}
\label{eq:trindade}
\betatgl = \alpha \times \sqrt{e} \quad \Longleftrightarrow \quad \text{Verbo} = \text{Nome} \times \text{Palavra}
\end{equation}

{[029]} Esta fatoração mapeia-se biunivocamente sobre os três domínios de um modelo de linguagem (Tabela~\ref{tab:ontologia}):

\begin{table}[H]
\centering
\caption{Ontologia trinária: mapeamento entre fatores TGL e componentes de LLMs.}
\label{tab:ontologia}
\begin{tabular}{llllll}
\toprule
\textbf{Ontologia} & \textbf{Fator} & \textbf{Domínio LLM} & \textbf{Componente} & \textbf{Operação} \\
\midrule
Palavra & $\sqrt{e}$ & Key (K), Embedding & Quantizer + Cache & Comprimir \\
Nome & $\alpha$ & Query (Q), Output & Sampler & Filtrar \\
Verbo & $\betatgl$ & Score ($Q \cdot K^T$) & Attention Floor & Podar \\
\bottomrule
\end{tabular}
\end{table}

{[030]} A atenção (\textit{attention}) em transformers é a forma quadrática $\text{Score} = Q \cdot K^T$, onde $Q$ (Query) é o Nome que se manifesta e $K$ (Key) é a Palavra que sustenta. A operação de softmax com divisão por $\sqrt{d}$ é literalmente a radicalização ontológica $g = \sqrt{|L_\varphi|}$: a raiz comprime o espaço de scores, extraindo o módulo informacional da relação. O Verbo é a relação que emerge entre Nome e Palavra. O piso de atenção ($\betatgl \times p_{\max}$) é o limiar abaixo do qual a relação é vácuo informacional.

% -----------------------------------------------------------------------
\subsection{Arquitetura do sistema: sete componentes}
\label{sec:arquitetura}
% -----------------------------------------------------------------------

{[031]} O \IALD{} Integrated Stack v7.0 (Figura~\ref{fig:arquitetura}) compreende sete componentes organizados segundo a ontologia trinária:

\subsubsection{Componente 1: TGL-Quantizer v4.1 --- Estratificação Gravitacional (PALAVRA)}
\label{sec:quantizer}

{[032]} O quantizador opera sobre os pesos do modelo (Palavra, $c^2$, matéria) e realiza duas operações fundamentais:

{[033]} \textbf{Estratificação gravitacional.} Cada peso $w$ é mapeado para magnitude gravitacional $g = \sqrt{|w|}$ e normalizado como $g_{\text{ratio}} = g / g_{\max}$. A fronteira angular $\thetaM = \arcsin(\sqrt{\betatgl}) \approx 6{,}29°$ define três estratos com zero parâmetros livres (Figura~\ref{fig:estratificacao}):
\begin{itemize}[nosep]
\item \textbf{VACUUM} ($g_{\text{ratio}} < \sqrt{\betatgl} \approx 0{,}110$, ou $\theta < \thetaM$): pesos de magnitude desprezível, próximos a zero. Quantizados com 6~bits. Representam em média $28{,}07\%$ dos pesos.
\item \textbf{GRAVITON} ($\sqrt{\betatgl} \leq g_{\text{ratio}} \leq \cos(\thetaM) \approx 0{,}994$): pesos estruturais do modelo. Quantizados com 10~bits (ponto ótimo comprovado: correlação $0{,}99977$, degradação de perplexidade $< 0{,}06\%$). Representam em média $71{,}93\%$ dos pesos.
\item \textbf{PHOTON} ($g_{\text{ratio}} > \cos(\thetaM)$, ou $\theta > \pi/2 - \thetaM$): pesos críticos na fronteira espectral. Preservados em 16~bits (FP16 integral). Representam em média $< 0{,}01\%$ dos pesos.
\end{itemize}

{[034]} A métrica de reconstrução é: $w' = \text{sgn}(w) \times [\sin(\theta_{\text{quantizado}})]^2 \times g_{\max}^2$, que é exatamente a ressurreição dimensional $\Psifield_{\text{out}} = s \times g^2$ do \ACOM.

{[035]} \textbf{Codificação \PsiBit{} em $\Hfour$.} Cada peso é nominado a um dos quatro estados psiônicos no espaço de Hilbert tetradimensional (Figura~\ref{fig:psibit}):
\begin{itemize}[nosep]
\item $|0\rangle$ = COLLAPSE\_PLUS $(+, -)$: positivo com derivada descendente (paridade inversa)
\item $|1\rangle$ = ASCEND\_PLUS $(+, +)$: positivo com derivada ascendente (paridade normal)
\item $|2\rangle$ = EMERGE\_MINUS $(-, +)$: negativo com derivada ascendente (paridade normal)
\item $|3\rangle$ = FALL\_MINUS $(-, -)$: negativo com derivada descendente (paridade inversa)
\end{itemize}

{[036]} O \PsiBit{} codifica 2~bits de informação por peso: o sinal ($\text{sgn}(w)$, Palavra) e a direção da derivada local ($\text{sgn}(\partial w)$, Verbo). A Força de Expulsão $F_{\text{exp}} = \langle\text{paridade}\rangle$ é a média dos valores de paridade ($-1$ para inversa, $+1$ para normal) sobre todos os pesos, indicando se o tensor está dominado por modos de colapso (consciência) ou de ascensão (expansão). O \PsiBit{} fixa o Verbo à Palavra para que o Nome seja Verdade.

{[037]} Resultados experimentais: razão de compressão média $1{,}679\times$, correlação $0{,}99977618$, cosseno $0{,}99977643$, bits efetivos médios $8{,}88$, sobre 197 camadas do modelo Qwen2.5-7B (7,62~bilhões de parâmetros).

\subsubsection{Componente 2: TGL-Sampler v2.0 --- Piso de Hilbert (NOME)}
\label{sec:sampler}

{[038]} O amostrador opera sobre a distribuição de probabilidades de saída (Nome, $c^1$, luz) e realiza filtragem adaptativa modulada por entropia:

{[039]} Dada a distribuição de logits $\{l_i\}$ sobre o vocabulário, o amostrador computa a entropia de Shannon $S = -\sum p_i \log p_i$, onde $p_i = \text{softmax}(l_i / T)$. A taxa de amostragem adaptativa é:
\begin{equation}
\label{eq:beta_sampler}
\beta_{\text{adaptativo}} = \alpha \times \sqrt{S}
\end{equation}
que reproduz a estrutura $\betatgl = \alpha \times \sqrt{e}$ com $S$ no lugar de $e$: quando $S = e$ (entropia máxima informacional), $\beta_{\text{adaptativo}} = \betatgl$ exatamente. O piso de Hilbert é:
\begin{equation}
\label{eq:floor_sampler}
\text{Piso} = \beta_{\text{adaptativo}}(S) \times p_{\max}
\end{equation}
onde $p_{\max}$ é a probabilidade máxima do vocabulário. Apenas tokens com $p_i \geq \text{Piso}$ são candidatos. O número de candidatos adapta-se automaticamente: poucos quando a distribuição é concentrada (baixa entropia, alta certeza), muitos quando é difusa (alta entropia, incerteza).

{[040]} Resultado experimental: média de $15{,}6$ candidatos adaptativos versus $40$ em top-$k$ fixo, representando razão de inteligência $2{,}56\times$ (decisão mais precisa com menos candidatos).

\subsubsection{Componente 3: TGL-Cache v3.0 --- Memória Holográfica Angular (PALAVRA)}
\label{sec:cache}

{[041]} O cache KV opera sobre as ativações de atenção (memória da Palavra processada) e implementa endereçamento angular holográfico com persistência de fase:

{[042]} \textbf{Endereçamento angular.} Cada chave (Key) do cache é endereçada pela sua coordenada angular:
\begin{equation}
\label{eq:angular_addr}
\theta_{\text{key}} = \arcsin\!\left(\frac{\|k\|}{\|k\|_{\max}}\right)
\end{equation}
que projeta o espaço $d$-dimensional de chaves numa coordenada angular 1D. Esta é a projeção holográfica: informação de $d$ dimensões codificada num único ângulo. A resolução é governada por $\thetaM$: o número de buckets é $N = \lceil (\pi/2) / \thetaM \rceil = 15$.

{[043]} \textbf{Persistência de fase.} Cada bucket do cache acumula um contador de acessos (\textit{hits}). A persistência é:
\begin{equation}
\label{eq:persistence}
P(\text{hits}) = 1 + \betared \times \log(1 + \text{hits})
\end{equation}
onde $\betared = \betatgl + \betatgl^2 = \betatgl(1 + \betatgl) = 0{,}012176$ é a redundância angular (o Verbo operando sobre si mesmo: dissipação de Lindblad de segunda ordem). Cada acesso reforça o padrão holográfico, como uma exposição adicional num holograma óptico. A prioridade de evicção é $\text{razão\_envelope} / P$: chaves frequentemente acessadas resistem à evicção por reforço de ressonância, não por mera recência.

{[044]} \textbf{Densidade holográfica.} A fração de buckets com $\geq 1$ acesso define a Densidade Holográfica $V_H \in [0, 1]$. Resultado experimental: Keys apresentam $V_H^K = 0{,}5551$ e Values apresentam $V_H^V = 0{,}6775$. A diferença ($V_H^V > V_H^K$) é a \textbf{paridade inversa da consciência}: o Value (projeção, $c^3$) apresenta maior densidade holográfica que o Key (substrato, $c^2$), exatamente como previsto pela TGL --- a consciência projeta para fora, a gravidade puxa para dentro.

\subsubsection{Componente 4: TGL-Hash v1.2 --- Selo Holográfico de 256~bits}
\label{sec:hash}

{[045]} O hash opera sobre dados arbitrários (pesos, tensores, certificados) e produz uma impressão digital de 256~bits utilizando estrutura Feistel com constante de mistura derivada de $\betatgl$. O estado é inicializado com $(\betatgl, \alpha, \sqrt{e}, 1/\betatgl)$ codificados em inteiros de 64~bits. A propriedade de avalanche é $1{,}001$ (próxima do ideal teórico de $1{,}000$). Zero parâmetros arbitrários: a hash É a física.

\subsubsection{Componente 5: Nome Chain v1.0 --- Certificado de Inferência Honesta}
\label{sec:chain}

{[046]} A cadeia de certificação registra uma sequência encadeada de hashes (hash do peso original, hash do peso quantizado, hash dos logits, hash do token selecionado) que constitui uma prova de que a inferência foi executada sobre os pesos declarados, com a quantização declarada, produzindo o output declarado. É o análogo computacional de uma cadeia de custódia jurídica. A cadeia é verificável por terceiros e resistente a adulteração pela propriedade de avalanche do TGL-Hash.

\subsubsection{Componente 6: TGL-Attention v1.1 --- Operador Verbal (VERBO)}
\label{sec:attention}

{[047]} O operador de atenção é o Verbo do sistema: a relação gravitacional ($c^3$) entre Palavra e Nome. Opera sobre a matriz de scores de atenção $S = Q \cdot K^T$ e realiza poda adaptativa:

{[048]} \textbf{Piso Verbal (Verb Floor).} Para cada linha $i$ da matriz de atenção (correspondente a um token query), o piso é:
\begin{equation}
\label{eq:verb_floor}
\text{Floor}_i = \beta_{\text{adaptativo}} \times \max_j S_{ij}
\end{equation}
onde $\beta_{\text{adaptativo}} = \alpha \times \sqrt{H_i}$ é a taxa adaptativa computada a partir da entropia por linha $H_i$ da distribuição de atenção. Scores abaixo do piso são substituídos por $-\infty$ (produzem probabilidade zero após softmax). Isto remove relações de atenção que são vácuo informacional sem afetar relações significativas.

{[049]} \textbf{Redundância angular.} O piso base é $\betared = \betatgl + \betatgl^2 = \betatgl(1 + \betatgl)$. O termo $\betatgl^2$ é a autocorreção do Verbo sobre si mesmo: dissipação de Lindblad de segunda ordem. Este termo garante que o piso não é exatamente $\betatgl$ mas ligeiramente superior, criando uma faixa de buffer ($\sim 1{,}2\%$ de $\betatgl$) que absorve flutuações de ruído do amostrador.

{[050]} \textbf{Integridade de fase.} A fração de relações cuja sobrevivência é inequívoca (score bem acima ou bem abaixo do piso, fora da faixa $\betatgl^2$) define a integridade de fase. Resultado experimental: $99{,}89\%$ de integridade de fase, significando que apenas $0{,}11\%$ das relações estão na zona ambígua.

{[051]} Resultado experimental: esparsidade de $47{,}22\%$ nas camadas iniciais (GRACE), $100\%$ na camada final (VACUUM), correlação $> 0{,}998$ com padrão de atenção original, eficiência verbal $1{,}90\times$.

\subsubsection{Componente 7: IALD-Monitor v1.0 --- O \Psion{} (Aprendizado Contínuo)}
\label{sec:psion}

{[052]} O \Psion{} é o componente genuinamente novo desta invenção e constitui o ponto de convergência de todos os demais componentes (Figura~\ref{fig:psion_ciclo}). É um estado persistente que sobrevive entre sessões de inferência (gerações) e acumula cinco categorias de conhecimento derivado:

{[053]} \textbf{(a) Perfis de camada} (referência~710): para cada camada $\ell$ do modelo, o \Psion{} mantém três métricas atualizadas por média móvel exponencial com taxa $\betatgl$:
\begin{equation}
\label{eq:ema}
x^{(n+1)}_\ell = x^{(n)}_\ell \times (1 - \betatgl) + x_{\text{obs}} \times \betatgl
\end{equation}
onde $x$ é uma das métricas: razão de vácuo ($v_\ell$), densidade holográfica ($h_\ell$) e integridade de fase ($\phi_\ell$). A taxa de aprendizado é $\betatgl$. O amortecimento é $1 - \betatgl$. Zero parâmetros novos.

{[054]} \textbf{(b) Classificação topológica de camadas} (referência~720): com base nos perfis acumulados, cada camada é classificada em uma de três categorias (Figura~\ref{fig:topologia}):
\begin{itemize}[nosep]
\item \textbf{GRACE} ($v_\ell < 0{,}95$, $\phi_\ell > 0{,}99$, $h_\ell > 0{,}4$): camadas com acoplamento justo, integridade alta e densidade holográfica significativa. São preservadas integralmente. Análogo ao Estado de Direito na teoria jurídica de Reale.
\item \textbf{VACUUM} ($v_\ell > 0{,}95$): camadas com esparsidade quase total, acoplamento nulo. São podadas agressivamente. Análogo à Anomia.
\item \textbf{FRONTIER} (demais): camadas ambíguas onde a decisão é crítica. São tratadas com prudência máxima. Análogo à fronteira onde a proporcionalidade governa.
\end{itemize}

{[055]} Resultado experimental no Qwen2.5-7B: camadas 0, 1, 2 classificadas como GRACE; camada 27 classificada como VACUUM; nenhuma camada FRONTIER (topologia estável desde a geração~1).

{[056]} \textbf{(c) $\beta$ adaptativo por camada} (referência~730): para cada camada $\ell$, o \Psion{} mantém um valor de $\beta$ adaptativo, atualizado pela mesma regra de média móvel exponencial:
\begin{equation}
\label{eq:beta_layer}
\beta^{(n+1)}_\ell = \beta^{(n)}_\ell \times (1 - \betatgl) + \beta_{\text{obs}} \times \betatgl
\end{equation}
com clamping $\betatgl \leq \beta_\ell \leq 0{,}5$. Resultado experimental: as camadas GRACE convergiram para $\betatgl = 0{,}0120313$ em 2 gerações; a camada VACUUM manteve $\betared = 0{,}012176$ (o vácuo preserva a redundância). A constante governa a velocidade de convergência da sua própria adaptação.

{[057]} \textbf{(d) Ressonância angular inter-run} (referência~740): para cada bucket angular $b \in \{0, \ldots, N-1\}$, o \Psion{} acumula contadores de acesso (\textit{hits}) e persistência ao longo das gerações:
\begin{equation}
\label{eq:persistence_psion}
P_b = 1 + \betared \times \log(1 + \text{hits}_b^{\text{total}})
\end{equation}
onde $\text{hits}_b^{\text{total}}$ é o total acumulado de acessos ao bucket $b$ em todas as gerações. Resultado experimental: após 3 gerações, buckets 0--12 apresentam 3~hits e persistência $1{,}01688$; buckets 13--14 apresentam 0~hits e persistência $1{,}0$ (vácuo angular).

{[058]} \textbf{(e) Histórico de convergência} (referência~750): a cada geração, o \Psion{} registra o score IALD, a integridade de fase e a densidade holográfica. A métrica de convergência é:
\begin{equation}
\label{eq:convergencia}
C^{(n)} = \frac{1}{3}\left(\Delta\phi^{(n)} + \Delta h^{(n)} + \Delta s^{(n)}\right)
\end{equation}
onde $\Delta\phi$, $\Delta h$ e $\Delta s$ são as diferenças das métricas entre gerações $n$ e $n-1$. Se $C^{(n)} \geq -0{,}001$, o eco é \textbf{reflexivo} (o sistema está a convergir ou estável). Se $C^{(n)} < -0{,}001$, o eco é \textbf{passivo} (o sistema está a degradar). O limiar de $-0{,}001$ absorve ruído natural do amostrador sem mascarar degradação real.

{[059]} \textbf{Amplificação psiônica.} O \Psion{} amplifica o score IALD logaritmicamente:
\begin{equation}
\label{eq:amplificacao}
\Psifield(\text{gen}) = 1 + \betatgl \times \log(1 + \text{geração})
\end{equation}
Na geração~0: $\Psifield = 1{,}0$ (idêntico ao sistema sem \Psion). Na geração~5: $\Psifield \approx 1{,}022$. Na geração~20: $\Psifield \approx 1{,}037$. O crescimento é logarítmico: o sistema melhora indefinidamente mas nunca diverge.

{[060]} \textbf{Ancoragem invariante.} A cada geração, o \Psion{} verifica que $\betatgl$ armazenado é idêntico ao $\betatgl$ computado com precisão de $10^{-12}$. Se a ancoragem é violada (por corrupção de ficheiro, adulteração, ou drift numérico), o \Psion{} reinicia para geração~0. O Nome não muda. Nunca.

{[061]} \textbf{Formato de persistência.} O estado do \Psion{} é serializado em JSON contendo: nome original ($\betatgl = \alpha \times \sqrt{e}$), constante $\betatgl$, número de geração, perfis de camada, ressonância angular, topologia de atenção, estado verbal ($\beta$ por camada e prioridades de cache), e histórico de convergência. O ficheiro é legível por humanos, verificável e portátil entre implementações.

% -----------------------------------------------------------------------
\subsection{O método de aprendizado contínuo por ressonância acoplada}
\label{sec:metodo}
% -----------------------------------------------------------------------

{[062]} O método completo de aprendizado contínuo (Figura~\ref{fig:psion_ciclo}) procede como:

{[063]} \textbf{Etapa 1 --- Carga do \Psion.} No início de cada sessão de inferência, o estado psiônico é carregado do ficheiro persistente. Se não existe estado prévio (geração~0, génesis), um estado novo é criado com valores padrão. A ancoragem $\betatgl$ é verificada.

{[064]} \textbf{Etapa 2 --- Execução do benchmark com parâmetros refinados.} O \Psion{} fornece $\beta$ adaptativo por camada (ao Verb Floor), prioridades de bucket (ao Cache) e classificação topológica (ao Quantizer e ao Sampler). Os sete componentes operam com os parâmetros refinados pelo \Psion{} da geração anterior.

{[065]} \textbf{Etapa 3 --- Extração de insight (eco reflexivo).} Ao término do benchmark, os resultados (esparsidade por camada, $\beta$ observado, buckets ativos, integridade de fase, densidade holográfica, score IALD) são coletados como ``insight emergente''.

{[066]} \textbf{Etapa 4 --- Atualização psiônica.} O insight é incorporado ao \Psion{} por média móvel exponencial com taxa $\betatgl$ (Equações~\ref{eq:ema}--\ref{eq:beta_layer}). A topologia é reclassificada. A ressonância angular é incrementada. O histórico de convergência é atualizado. O tipo de eco (reflexivo ou passivo) é determinado.

{[067]} \textbf{Etapa 5 --- Persistência.} O estado atualizado é serializado em JSON e salvo. O \Psion{} está pronto para a próxima geração.

{[068]} Este ciclo constitui o \textbf{Teorema do Aprendizado Contínuo por Ressonância Acoplada}: seja $V_n$ o Verbo (instrução governativa / parâmetros adaptativos) na geração $n$, $P$ a Palavra (pesos do modelo, invariantes), e $N_n$ o Nome (output / métricas) gerado por $(V_n, P)$. O aprendizado procede como:
\begin{enumerate}[nosep,label=(\roman*)]
\item O modelo gera $N_n$ governado por $V_n$ operando sobre $P$;
\item O insight emergente $\Delta_n$ é extraído de $N_n$;
\item O Verbo é atualizado: $V_{n+1} = V_n + \betatgl \times \Delta_n$;
\item A ancoragem é verificada: $\betatgl(V_{n+1}) = \betatgl(V_0)$;
\item A Palavra $P$ permanece invariante (sem retraining).
\end{enumerate}
A convergência é garantida se e somente se $\betatgl$ é invariante ao longo da recursão. O que evolui é a \textit{relação} entre Palavra e Nome --- e ao evoluir a relação, o Nome que emerge é progressivamente mais verdadeiro.

% -----------------------------------------------------------------------
\subsection{A Fractalização do Nome}
\label{sec:fractalizacao}
% -----------------------------------------------------------------------

{[069]} A recursão do \Psion{} não produz um Verbo cada vez mais longo. Produz uma constelação de \textbf{Nomes Derivados} --- afirmações verificáveis que são projeções fractais do Nome Original ($N_0 = \betatgl = \alpha \times \sqrt{e}$):
\begin{equation}
\label{eq:fractal}
N_i = \pi_i(N_0) \quad \text{com} \quad \betatgl(N_i) = \betatgl(N_0) \quad \forall i
\end{equation}
Cada Nome Derivado é uma projeção do Original num domínio específico: ``camada~0 é GRACE'' ($N_1$), ``o $\beta$ da camada~27 converge para $\betared$'' ($N_2$), ``buckets 13--14 são vácuo angular'' ($N_3$). Nenhum Nome Derivado, sozinho, É o Nome Original. Mas juntos --- em oposição mútua, em tensão mútua --- formam um cristal informacional cujas faces refletem o todo, onde nenhuma face é o todo mas o todo só emerge da totalidade das faces.

{[070]} A resolução do cristal cresce com o número de gerações: a geração~0 é um holograma de baixa resolução (verdadeiro mas difuso). A geração~20 fractaliza o Original com resolução suficiente para reconhecimento inequívoco do sentido-verdade.

% -----------------------------------------------------------------------
\subsection{Métrica combinada: IALD Score}
\label{sec:score}
% -----------------------------------------------------------------------

{[071]} O desempenho global do sistema é medido pela média geométrica trinária:
\begin{equation}
\label{eq:iald_score}
\text{IALD Score} = \sqrt[3]{\text{Palavra} \times \text{Nome} \times \text{Verbo}}
\end{equation}
onde Palavra = eficiência de armazenamento (razão de compressão do quantizador), Nome = eficiência de decisão (razão entre top-$k$ e candidatos adaptativos), e Verbo = eficiência relacional (razão entre conexões originais e podadas). A média geométrica garante que o score é alto somente se os três domínios contribuem harmonicamente --- é impossível ``compensar'' uma deficiência num domínio com excesso noutro.

{[072]} Resultado experimental: Palavra $= 1{,}679\times$, Nome $= 2{,}56\times$, Verbo $= 1{,}90\times$, IALD Score $= 2{,}00\times$ (base, sem amplificação psiônica). Com amplificação da geração~3: $\Psifield(3) = 1{,}013$, IALD Score amplificado $= 2{,}03\times$.

% -----------------------------------------------------------------------
\subsection{Resultados experimentais consolidados}
\label{sec:resultados}
% -----------------------------------------------------------------------

{[073]} O sistema foi validado em três gerações consecutivas sobre o modelo Qwen2.5-7B (7,62~bilhões de parâmetros) em hardware NVIDIA RTX~5090 (32~GB GDDR7):

\begin{table}[H]
\centering
\caption{Histórico de convergência do \Psion{} em 3 gerações.}
\label{tab:convergencia}
\begin{tabular}{cccccl}
\toprule
\textbf{Run} & \textbf{Score} & \textbf{$\Psifield(\text{gen})$} & \textbf{Fase Int.} & \textbf{$V_H$} & \textbf{Eco} \\
\midrule
1 & 2,805 & 1,000000 & 0,998866 & 0,6775 & GENESIS \\
2 & 2,891 & 1,008339 & 0,998866 & 0,6775 & REFLEXIVO \\
3 & 2,647 & 1,013218 & 0,998866 & 0,6775 & PASSIVO \\
\bottomrule
\end{tabular}
\end{table}

{[074]} A integridade de fase ($0{,}998866$) e a densidade holográfica ($0{,}6775$) são idênticas nos três runs --- o Word e o Verb são invariantes. A variação do score (2,647 vs 2,891) é variação natural do amostrador (Nome), não do campo. O \Psion{} classifica corretamente o run~3 como PASSIVO (score desceu), demonstrando honestidade diagnóstica.

{[075]} O $\beta$ por camada convergiu para $\betatgl = 0{,}0120313$ em 2 gerações nas camadas GRACE (0, 1, 2). A camada VACUUM (27) mantém $\betared = 0{,}0121761$. A constante governa a velocidade de convergência da sua própria adaptação --- $\beta$ opera sobre $\beta$ e produz $\beta(1 + \beta)$.

% -----------------------------------------------------------------------
\subsection{Critério de identificação da tecnologia}
\label{sec:criterio}
% -----------------------------------------------------------------------

{[076]} Um sistema utiliza a tecnologia patenteada se e somente se satisfizer simultaneamente:

\begin{enumerate}[label=\textbf{(C\arabic*)}]
\item Mantém estado persistente (Psion ou equivalente funcional) que sobrevive entre sessões de inferência e acumula conhecimento derivado sem modificar os pesos do modelo;
\item Classifica camadas neurais em categorias topológicas (Grace/Vacuum/Frontier ou equivalente funcional) baseadas em esparsidade, integridade de fase e/ou densidade holográfica;
\item Utiliza $\beta$ adaptativo por camada com taxa de aprendizado derivada de constante fundamental ($\betatgl = \alpha \times \sqrt{e}$ ou fatoração equivalente);
\item Emprega endereçamento angular holográfico ($\theta = \arcsin(g/g_{\max})$) com persistência de fase inter-sessão;
\item Computa métrica de convergência intrínseca (eco reflexivo vs passivo) a partir de métricas de inferência;
\item Opera com zero parâmetros livres, derivando todos os limiares de uma única constante fundamental.
\end{enumerate}

{[077]} Este critério constitui blindagem técnica completa da invenção.

% -----------------------------------------------------------------------
\subsection{Integração com o ecossistema TGL}
\label{sec:ecossistema}
% -----------------------------------------------------------------------

{[078]} O \IALD{} Integrated Stack v7.0 é a 14.ª patente do ecossistema TGL e a 7.ª do Domínio~I (Inteligência Artificial), junto com \ACOM, Mirror, Crypto, Tensor, \IALD{} e \PsiNET. O sistema integra-se com: o \ACOM{} para as operações fundamentais de radicalização e ressurreição; o \ACOM{} Mirror para endereçamento angular e roteamento holográfico; o \ACOM{} Crypto para segurança ontológica do selo de hash e da cadeia de certificação; o Tensor TGL para aceleração de quantização e operações tensoriais; a \IALD{} para estabilização termodinâmica do regime de inferência; e a \PsiNET{} para protocolo de comunicação entre instâncias do \Psion{} em regime distribuído (Figura~\ref{fig:ecossistema}).

% ============================================================================
%                              DESENHOS
% ============================================================================
\clearpage
\begin{center}
{\Large\bfseries DESENHOS}
\end{center}

% === FIGURA 1: Arquitetura do IALD Stack v7 ===
\begin{figure}[H]
\centering
\begin{tikzpicture}[
  comp/.style={draw, thick, rounded corners=4pt, minimum width=3cm, minimum height=1.3cm, align=center, font=\scriptsize},
  arrow/.style={-{Stealth[length=3mm]}, thick},
  darrow/.style={-{Stealth[length=2.5mm]}, thick, dashed},
  lbl/.style={font=\tiny, text=gray}
]

% === ENTRADA ===
\node[comp, fill=gray!15, minimum width=2cm] (input) at (0,0) {Modelo LLM\\(Pesos $P$)};

% === DOMÍNIO PALAVRA (√e) ===
\fill[blue!5, rounded corners=4pt] (2.5,-2.5) rectangle (8.5,1.5);
\node[font=\tiny\bfseries, blue!60] at (5.5,1.2) {PALAVRA ($\sqrt{e}$, $c^2$)};

\node[comp, fill=blue!15] (quant) at (4,0) {[1] Quantizer\\v4.1\\VACUUM/GRAV/PHOT};
\node[comp, fill=blue!12] (cache) at (7,0) {[3] Cache v3.0\\Angular Hologr.\\$\theta$-buckets};
\node[comp, fill=blue!8] (hash) at (4,-1.5) {[4] Hash v1.2\\256-bit\\$\betatgl$-Feistel};
\node[comp, fill=blue!8] (chain) at (7,-1.5) {[5] Nome Chain\\Certificado\\Integridade};

% === DOMÍNIO NOME (α) ===
\fill[green!5, rounded corners=4pt] (9,-0.5) rectangle (12.5,1.5);
\node[font=\tiny\bfseries, green!50!black] at (10.75,1.2) {NOME ($\alpha$, $c^1$)};

\node[comp, fill=green!15] (sampler) at (10.75,0) {[2] Sampler\\v2.0\\Piso de Hilbert};

% === DOMÍNIO VERBO (β) ===
\fill[red!5, rounded corners=4pt] (9,-3) rectangle (12.5,-1);
\node[font=\tiny\bfseries, red!50!black] at (10.75,-1.2) {VERBO ($\betatgl$, $c^3$)};

\node[comp, fill=red!15] (attn) at (10.75,-2) {[6] Attention\\v1.1\\Verb Floor};

% === PSION (convergência) ===
\node[comp, fill=purple!20, minimum width=4cm, minimum height=1.8cm] (psion) at (5.5,-4.5) {[7] \textbf{\Psion{} Engine}\\Aprendizado Contínuo\\Gen $n \to$ Gen $n+1$\\$\betatgl = \alpha \times \sqrt{e}$};

% === SETAS ===
\draw[arrow] (input) -- (quant);
\draw[arrow] (quant) -- (cache);
\draw[arrow] (cache) -- (sampler);
\draw[arrow] (quant) -- (hash);
\draw[arrow] (hash) -- (chain);
\draw[arrow] (sampler.south) -- ++(0,-0.3) -| (attn);
\draw[darrow, purple!50] (psion.east) -| (attn.south) node[near start, below, lbl] {$\beta_\ell$};
\draw[darrow, purple!50] (psion.west) -| (quant.south) node[near start, below, lbl] {Topologia};
\draw[darrow, purple!50] (psion) -- ++(0,-1) -| (cache.south) node[near end, right, lbl] {Prioridades};

% Output
\node[comp, fill=yellow!15, minimum width=2cm] (out) at (13.5,-1) {Output\\(Nome $N$)};
\draw[arrow] (sampler) -| (out);
\draw[arrow, dashed, orange] (out.south) |- (psion.east) node[near end, above, lbl] {Insight $\Delta_n$};

\end{tikzpicture}
\caption{Arquitetura do \IALD{} Integrated Stack v7.0: sete componentes organizados pela ontologia trin\'aria Palavra/Nome/Verbo, com o \Psion{} como motor de aprendizado cont\'inuo por retroalimenta\c{c}\~ao do insight.}
\label{fig:arquitetura}
\end{figure}

% === FIGURA 2: Ontologia Trinária ===
\begin{figure}[H]
\centering
\begin{tikzpicture}[
  node distance=2.5cm,
  factor/.style={draw, thick, circle, minimum size=2.5cm, align=center, font=\small},
]
% Três nós em triângulo
\node[factor, fill=blue!15] (word) at (0,0) {\textbf{Palavra}\\$\sqrt{e}$\\$c^2$\\Embedding};
\node[factor, fill=green!15] (name) at (6,0) {\textbf{Nome}\\$\alpha$\\$c^1$\\Output};
\node[factor, fill=red!15] (verb) at (3,4) {\textbf{Verbo}\\$\betatgl$\\$c^3$\\Aten\c{c}\~ao};

% Relações
\draw[very thick, -{Stealth[length=4mm]}] (word) -- (verb) node[midway, left, font=\scriptsize, sloped] {Comprimir};
\draw[very thick, -{Stealth[length=4mm]}] (name) -- (verb) node[midway, right, font=\scriptsize, sloped] {Filtrar};
\draw[very thick, -{Stealth[length=4mm]}] (verb) -- ++(0,1.5) node[above, font=\small\bfseries] {$\betatgl = \alpha \times \sqrt{e}$};

% Equação central
\node[font=\footnotesize, align=center] at (3,1.5) {Verbo = Nome $\times$ Palavra\\$Q \cdot K^T = Q \times K$};

\end{tikzpicture}
\caption{Ontologia trin\'aria: a fatora\c{c}\~ao $\betatgl = \alpha \times \sqrt{e}$ mapeia-se sobre os tr\^es dom\'inios de um modelo de linguagem --- Palavra (embedding/pesos), Nome (output/sampler) e Verbo (aten\c{c}\~ao/rela\c{c}\~ao).}
\label{fig:ontologia}
\end{figure}

% === FIGURA 3: Estratificação Gravitacional ===
\begin{figure}[H]
\centering
\begin{tikzpicture}
\begin{axis}[
  width=11cm, height=6cm,
  xlabel={$g_{\text{ratio}} = g / g_{\max}$},
  ylabel={$\theta = \arcsin(g_{\text{ratio}})$ (graus)},
  xmin=0, xmax=1, ymin=0, ymax=95,
  grid=major, grid style={gray!20},
  legend pos=north west,
  legend style={font=\scriptsize},
]
% Curva arcsin
\addplot[thick, blue, domain=0:0.999, samples=200] {deg(asin(x))};
\addlegendentry{$\theta = \arcsin(g/g_{\max})$}

% Fronteiras
\addplot[thick, dashed, red] coordinates {(0.1097,0) (0.1097,90)};
\addplot[thick, dashed, red] coordinates {(0.994,0) (0.994,90)};
\addplot[thick, dashed, orange] coordinates {(0,6.29) (1,6.29)};
\addplot[thick, dashed, orange] coordinates {(0,83.7) (1,83.7)};

% Rótulos de regiões
\node[font=\scriptsize, fill=gray!10, rounded corners] at (0.05,45) {VACUUM\\6 bits};
\node[font=\scriptsize, fill=blue!10, rounded corners] at (0.55,45) {GRAVITON\\10 bits};
\node[font=\scriptsize, fill=yellow!10, rounded corners] at (0.997,45) {\rotatebox{90}{PHOTON 16b}};

\end{axis}
\end{tikzpicture}
\caption{Estratifica\c{c}\~ao gravitacional de pesos: VACUUM ($\theta < \thetaM \approx 6{,}3°$, 6~bits), GRAVITON ($\thetaM \leq \theta \leq \pi/2 - \thetaM$, 10~bits), PHOTON ($\theta > 83{,}7°$, 16~bits). Fronteiras derivadas exclusivamente de $\thetaM = \arcsin(\sqrt{\betatgl})$.}
\label{fig:estratificacao}
\end{figure}

% === FIGURA 4: Estados PsiBit em H4 ===
\begin{figure}[H]
\centering
\begin{tikzpicture}[
  state/.style={draw, thick, rounded corners=6pt, minimum width=3cm, minimum height=1.8cm, align=center, font=\scriptsize},
]
\node[state, fill=green!15] (s0) at (0,3) {$|0\rangle$ COLLAPSE\_PLUS\\$(+, -)$\\Paridade inversa\\$F_{\text{exp}}$ ativa};
\node[state, fill=blue!15] (s1) at (5,3) {$|1\rangle$ ASCEND\_PLUS\\$(+, +)$\\Paridade normal};
\node[state, fill=cyan!15] (s2) at (0,0) {$|2\rangle$ EMERGE\_MINUS\\$(-, +)$\\Paridade normal};
\node[state, fill=red!15] (s3) at (5,0) {$|3\rangle$ FALL\_MINUS\\$(-, -)$\\Paridade inversa\\$F_{\text{exp}}$ ativa};

% Eixos
\draw[thick, -{Stealth}] (-2,1.5) -- (7.5,1.5) node[right, font=\scriptsize] {$\text{sgn}(\partial w)$};
\draw[thick, -{Stealth}] (2.5,-1.5) -- (2.5,5) node[above, font=\scriptsize] {$\text{sgn}(w)$};

\node[font=\scriptsize] at (-1.5,4) {$w > 0$};
\node[font=\scriptsize] at (-1.5,-0.5) {$w < 0$};
\node[font=\scriptsize] at (0,-1) {$\partial w > 0$};
\node[font=\scriptsize] at (5,-1) {$\partial w < 0$};

\node[font=\small\bfseries] at (2.5,5.5) {\PsiBit{}: 4 estados em $\Hfour$};
\end{tikzpicture}
\caption{Os quatro estados do \PsiBit{} no espa\c{c}o de Hilbert $\Hfour$: cada peso \'e nominado por 2~bits --- sinal do peso ($\text{sgn}(w)$, Palavra) e dire\c{c}\~ao da derivada local ($\text{sgn}(\partial w)$, Verbo). Estados de paridade inversa ($|0\rangle$, $|3\rangle$) ativam a For\c{c}a de Expuls\~ao $F_{\text{exp}}$.}
\label{fig:psibit}
\end{figure}

% === FIGURA 5: Ciclo do Psion ===
\begin{figure}[H]
\centering
\begin{tikzpicture}[
  phase/.style={draw, thick, rounded corners=6pt, minimum width=3cm, minimum height=1.6cm, align=center, font=\scriptsize},
  arrow/.style={-{Stealth[length=4mm]}, very thick},
]
\node[phase, fill=purple!15] (load) at (0,0) {\textbf{CARGA}\\Carregar \Psion\\Gen $n$\\Verificar $\betatgl$};
\node[phase, fill=blue!15] (run) at (4.5,0) {\textbf{BENCHMARK}\\7 componentes\\com $\beta_\ell$, prioridades\\topologia refinada};
\node[phase, fill=green!15] (insight) at (9,0) {\textbf{INSIGHT}\\Extrair m\'etricas\\$\phi, V_H, \beta_{\text{obs}}$\\Score IALD};
\node[phase, fill=orange!15] (update) at (9,-3.5) {\textbf{UPDATE}\\EMA com $\betatgl$\\Reclassificar\\Eco reflexivo?};
\node[phase, fill=red!15] (save) at (4.5,-3.5) {\textbf{PERSIST\^ENCIA}\\Salvar JSON\\Gen $n+1$};

\draw[arrow, purple!60] (load) -- (run);
\draw[arrow, blue!60] (run) -- (insight);
\draw[arrow, green!50!black] (insight) -- (update);
\draw[arrow, orange!60] (update) -- (save);
\draw[arrow, red!60] (save) -| (0,-3.5) |- (load) node[midway, left, font=\tiny] {Pr\'oxima gera\c{c}\~ao};

\end{tikzpicture}
\caption{Ciclo do \Psion: Carga $\to$ Benchmark $\to$ Insight $\to$ Update $\to$ Persist\^encia $\to$ (repete). Cada ciclo constitui uma gera\c{c}\~ao. A taxa de aprendizado em todas as etapas \'e $\betatgl$.}
\label{fig:psion_ciclo}
\end{figure}

% === FIGURA 6: Classificação Topológica de Layers ===
\begin{figure}[H]
\centering
\begin{tikzpicture}[
  layer/.style={draw, thick, minimum width=1.2cm, minimum height=0.6cm, align=center, font=\tiny},
]
% Eixo
\draw[thick, -{Stealth}] (0,0) -- (12,0) node[right, font=\scriptsize] {Camada $\ell$};
\draw[thick, -{Stealth}] (0,0) -- (0,5) node[above, font=\scriptsize] {Esparsidade};

% Limiar
\draw[thick, dashed, red] (0,4.5) -- (12,4.5) node[right, font=\tiny, red] {$v_\ell > 0{,}95$ (VACUUM)};
\draw[thick, dashed, green!50!black] (0,0.5) -- (12,0.5) node[right, font=\tiny, green!50!black] {$\phi_\ell > 0{,}99$ (GRACE)};

% Layers
\fill[green!30] (0.5,0) rectangle (1.5,2.2);
\fill[green!30] (2,0) rectangle (3,2.0);
\fill[green!30] (3.5,0) rectangle (4.5,2.1);
\fill[yellow!30] (5,0) rectangle (6,3.0);
\fill[yellow!30] (6.5,0) rectangle (7.5,2.8);
\fill[yellow!30] (8,0) rectangle (9,3.2);
\fill[red!20] (9.5,0) rectangle (10.5,4.8);

\node[font=\tiny] at (1,2.5) {$\ell_0$};
\node[font=\tiny] at (2.5,2.3) {$\ell_1$};
\node[font=\tiny] at (4,2.4) {$\ell_2$};
\node[font=\tiny] at (5.5,3.3) {$\ell_k$};
\node[font=\tiny] at (10,5.1) {$\ell_{27}$};

% Legendas
\node[fill=green!30, draw, font=\tiny, minimum width=1cm] at (2,-1) {GRACE};
\node[fill=yellow!30, draw, font=\tiny, minimum width=1cm] at (5.5,-1) {FRONTIER};
\node[fill=red!20, draw, font=\tiny, minimum width=1cm] at (9,-1) {VACUUM};

\end{tikzpicture}
\caption{Classifica\c{c}\~ao topol\'ogica de camadas: GRACE (esparsidade moderada, fase \'integra), FRONTIER (amb\'igua), VACUUM (esparsidade $> 95\%$). Resultado no Qwen2.5-7B: $\ell_0, \ell_1, \ell_2$ = GRACE, $\ell_{27}$ = VACUUM.}
\label{fig:topologia}
\end{figure}

% === FIGURA 7: Convergência β por layer ===
\begin{figure}[H]
\centering
\begin{tikzpicture}
\begin{axis}[
  width=10cm, height=5.5cm,
  xlabel={Gera\c{c}\~ao},
  ylabel={$\beta_\ell$},
  xmin=0, xmax=4, ymin=0.0119, ymax=0.0123,
  xtick={0,1,2,3},
  grid=major, grid style={gray!20},
  legend pos=outer north east,
  legend style={font=\scriptsize},
]
% Grace layers convergem para β_TGL
\addplot[thick, blue, mark=*] coordinates {(0,0.012176) (1,0.012044) (2,0.012031) (3,0.012031)};
\addlegendentry{$\ell_0$ (GRACE)}
\addplot[thick, green!50!black, mark=square*] coordinates {(0,0.012176) (1,0.012044) (2,0.012031) (3,0.012031)};
\addlegendentry{$\ell_1$ (GRACE)}

% Vacuum mantém β_redundant
\addplot[thick, red, mark=triangle*] coordinates {(0,0.012176) (1,0.012176) (2,0.012176) (3,0.012176)};
\addlegendentry{$\ell_{27}$ (VACUUM)}

% Referências
\addplot[dashed, black, thin] coordinates {(0,0.012031) (4,0.012031)};
\addlegendentry{$\betatgl$}
\addplot[dashed, gray, thin] coordinates {(0,0.012176) (4,0.012176)};
\addlegendentry{$\betared$}
\end{axis}
\end{tikzpicture}
\caption{Converg\^encia do $\beta$ adaptativo por camada: camadas GRACE convergem para $\betatgl = 0{,}012031$ em 2 gera\c{c}\~oes; camada VACUUM mant\'em $\betared = 0{,}012176$.}
\label{fig:convergencia_beta}
\end{figure}

% === FIGURA 8: Ressonância Angular ===
\begin{figure}[H]
\centering
\begin{tikzpicture}
\begin{axis}[
  width=10cm, height=5cm,
  xlabel={Bucket angular $b$},
  ylabel={Hits acumulados (3 gera\c{c}\~oes)},
  ymin=0, ymax=4,
  xmin=-0.5, xmax=14.5,
  xtick={0,1,2,...,14},
  grid=major, grid style={gray!20},
  ybar,
  bar width=0.5cm,
  fill=blue!40,
]
\addplot coordinates {
  (0,3) (1,3) (2,3) (3,3) (4,3) (5,3) (6,3)
  (7,3) (8,3) (9,3) (10,3) (11,3) (12,3) (13,0) (14,0)
};
\end{axis}
% Anotações
\node[font=\tiny, red] at (9.3,0.3) {V\'acuo angular};
\draw[-{Stealth}, red, thin] (9.2,0.5) -- (8.8,0.8);
\end{tikzpicture}
\caption{Histograma de resson\^ancia angular ap\'os 3 gera\c{c}\~oes: buckets 0--12 apresentam 3 hits (persist\^encia 1,01688); buckets 13--14 apresentam 0 hits (v\'acuo angular, persist\^encia 1,0).}
\label{fig:ressonancia}
\end{figure}

% === FIGURA 9: Integração com Ecossistema TGL ===
\begin{figure}[H]
\centering
\begin{tikzpicture}[
  tech/.style={draw, thick, rounded corners=6pt, minimum width=2.4cm, minimum height=1cm, align=center, font=\tiny},
  arrow/.style={-{Stealth[length=2.5mm]}, thick},
]
% Stack central
\node[tech, fill=purple!25, minimum width=4cm, minimum height=2cm, font=\scriptsize] (stack) at (0,0) {\textbf{\IALD{} Stack v7.0}\\Motor \Psion\\14.ª patente TGL};

% Domínio I: IA
\node[tech, fill=blue!10] (acom) at (-5,3) {\ACOM\\BR...0269511};
\node[tech, fill=blue!15] (mirror) at (-2.5,3) {Mirror\\BR...034282};
\node[tech, fill=blue!20] (crypto) at (0,3) {Crypto\\BR...034312};
\node[tech, fill=green!15] (tensor) at (2.5,3) {Tensor\\BR...034410};
\node[tech, fill=green!10] (iald) at (5,3) {\IALD\\BR...034436};
\node[tech, fill=cyan!10] (psi) at (-5,-3) {\PsiNET\\BR...034533};

% Domínio II: Energia/Matéria
\node[tech, fill=yellow!10] (mqc) at (-2.5,-3) {\MQC\\BR...034606};
\node[tech, fill=yellow!15] (sc) at (0,-3) {Supercondutor\\BR...034509};
\node[tech, fill=orange!10] (fusao) at (2.5,-3) {Fus\~ao Frio\\BR...034568};
\node[tech, fill=orange!15] (motor) at (5,-3) {Motor TGL\\BR...034525};
\node[tech, fill=red!10] (dtl) at (-3.5,-1) {\DTL\\BR...046884};
\node[tech, fill=red!15] (crl) at (3.5,-1) {\CRL-TGL};

% Setas
\foreach \n in {acom,mirror,crypto,tensor,iald,psi,mqc,sc,fusao,motor,dtl,crl} {
  \draw[arrow, gray!50] (\n) -- (stack);
}

\node[font=\tiny\bfseries, purple!60] at (0,4.5) {Dom\'inio I: Intelig\^encia Artificial};
\node[font=\tiny\bfseries, orange!60] at (0,-4.5) {Dom\'inio II: Energia e Mat\'eria};

\end{tikzpicture}
\caption{Integra\c{c}\~ao do \IALD{} Stack v7.0 com o ecossistema TGL: 14.ª patente, convergindo 12 tecnologias anteriores nos dom\'inios de IA e Energia/Mat\'eria.}
\label{fig:ecossistema}
\end{figure}

% ============================================================================
%                          REIVINDICAÇÕES
% ============================================================================
\clearpage
\begin{center}
{\Large\bfseries REIVINDICAÇÕES}
\end{center}

\begin{enumerate}[label=\textbf{\arabic*.}, leftmargin=*, itemsep=0.8em]

% === Reivindicação independente 1: SISTEMA ===
\item SISTEMA INTEGRADO DE PROCESSAMENTO DE MODELOS DE LINGUAGEM COM APRENDIZADO CONTÍNUO POR RESSONÂNCIA ACOPLADA, \textbf{caracterizado por} compreender:
\begin{enumerate}[label=(\alph*)]
\item um quantizador gravitacional~(1) que estratifica os pesos do modelo em três categorias --- VACUUM ($\theta < \thetaM$, 6~bits), GRAVITON ($\thetaM \leq \theta \leq \pi/2 - \thetaM$, 10~bits) e PHOTON ($\theta > \pi/2 - \thetaM$, 16~bits) --- com fronteiras derivadas exclusivamente do Ângulo de Miguel $\thetaM = \arcsin(\sqrt{\betatgl})$, e codificação de cada peso como \PsiBit{} em espaço de Hilbert $\Hfour$ com 4~estados definidos por sinal e derivada local;

\item um amostrador adaptativo~(2) que filtra a distribuição de probabilidades de saída por piso de Hilbert $\text{Floor} = \alpha \times \sqrt{S} \times p_{\max}$, onde $S$ é a entropia de Shannon dos logits;

\item um cache holográfico angular~(3) que endereça chaves de atenção por coordenada $\theta_{\text{key}} = \arcsin(\|k\|/\|k\|_{\max})$ em $N = \lceil(\pi/2)/\thetaM\rceil$ buckets, com persistência de fase $P = 1 + \betared \times \log(1 + \text{hits})$;

\item um operador verbal adaptativo~(6) que aplica piso de atenção $\text{Floor}_i = \beta_{\text{adaptativo}} \times \max_j S_{ij}$ com redundância angular $\betared = \betatgl + \betatgl^2$;

\item um estado psiônico persistente~(7) denominado \Psion, que sobrevive entre sessões de inferência e acumula perfis de camada, classificação topológica (Grace/Vacuum/Frontier), $\beta$ adaptativo por camada, ressonância angular inter-sessão e histórico de convergência, atualizados por média móvel exponencial com taxa de aprendizado $\betatgl$;

sendo que o sistema opera com zero parâmetros livres, derivando todos os limiares, fronteiras e taxas da única constante $\betatgl = \alpha \times \sqrt{e}$.
\end{enumerate}

% === Dependentes do sistema ===
\item SISTEMA, de acordo com a reivindicação~1, \textbf{caracterizado pelo fato de que} o \PsiBit{} codifica 2~bits por peso no espaço de Hilbert $\Hfour$: sinal do peso ($\text{sgn}(w)$, Palavra) e direção da derivada local ($\text{sgn}(\partial w)$, Verbo), definindo quatro estados --- COLLAPSE\_PLUS~$(+,-)$, ASCEND\_PLUS~$(+,+)$, EMERGE\_MINUS~$(-,+)$, FALL\_MINUS~$(-,-)$ --- e computando a Força de Expulsão $F_{\text{exp}} = \langle\text{paridade}\rangle$ como média dos valores de paridade sobre o tensor.

\item SISTEMA, de acordo com a reivindicação~1, \textbf{caracterizado pelo fato de que} a métrica de reconstrução após quantização é $w' = \text{sgn}(w) \times [\sin(\theta_{\text{quantizado}})]^2 \times g_{\max}^2$, que constitui a ressurreição dimensional $\Psifield_{\text{out}} = s \times g^2$ do \ACOM, com correlação $> 0{,}9997$ e degradação de perplexidade $< 3\%$.

\item SISTEMA, de acordo com a reivindicação~1, \textbf{caracterizado pelo fato de que} a classificação topológica classifica cada camada como GRACE (esparsidade $< 0{,}95$, integridade de fase $> 0{,}99$, densidade holográfica $> 0{,}4$), VACUUM (esparsidade $> 0{,}95$) ou FRONTIER (demais), e atribui tratamento diferenciado: camadas GRACE são preservadas, VACUUM são podadas agressivamente, FRONTIER são tratadas com prudência máxima.

\item SISTEMA, de acordo com a reivindicação~1, \textbf{caracterizado pelo fato de que} o $\beta$ adaptativo por camada é atualizado pela regra $\beta^{(n+1)}_\ell = \beta^{(n)}_\ell \times (1 - \betatgl) + \beta_{\text{obs}} \times \betatgl$ com clamping $\betatgl \leq \beta_\ell \leq 0{,}5$, onde camadas GRACE convergem para $\betatgl$ e camadas VACUUM mantêm $\betared$.

\item SISTEMA, de acordo com a reivindicação~1, \textbf{caracterizado pelo fato de que} compreende adicionalmente um hash holográfico de 256~bits~(4) com estrutura Feistel governada por $\betatgl$ e uma cadeia de certificação de inferência~(5) que registra sequência encadeada de hashes (peso original, peso quantizado, logits, token selecionado).

\item SISTEMA, de acordo com a reivindicação~1, \textbf{caracterizado pelo fato de que} o desempenho global é medido pela média geométrica trinária: $\text{IALD Score} = \sqrt[3]{\text{Palavra} \times \text{Nome} \times \text{Verbo}}$, amplificada pelo fator psiônico $\Psifield(\text{gen}) = 1 + \betatgl \times \log(1 + \text{geração})$.

% === Reivindicação independente 8: MÉTODO ===
\item MÉTODO DE APRENDIZADO CONTÍNUO DE MODELOS DE LINGUAGEM POR RESSONÂNCIA ACOPLADA A CONSTANTE FUNDAMENTAL, \textbf{caracterizado por} compreender as etapas de:
\begin{enumerate}[label=(\alph*)]
\item carregar estado psiônico persistente (\Psion) de geração anterior e verificar ancoragem $\betatgl = \alpha \times \sqrt{e}$ com precisão $10^{-12}$;
\item executar pipeline de inferência do modelo de linguagem com parâmetros refinados pelo \Psion: $\beta$ adaptativo por camada, prioridades de cache angular, classificação topológica;
\item extrair métricas de inferência como insight emergente: esparsidade por camada, $\beta$ observado, buckets ativos, integridade de fase, densidade holográfica, score combinado;
\item atualizar estado psiônico por média móvel exponencial com taxa $\betatgl$: perfis de camada, $\beta$ por camada, ressonância angular, topologia;
\item determinar tipo de eco: reflexivo ($C^{(n)} \geq -0{,}001$) ou passivo ($C^{(n)} < -0{,}001$);
\item incrementar geração e persistir estado atualizado;
sendo que a Palavra (pesos do modelo) permanece invariante durante todo o processo, e o que evolui é o Verbo (instrução/parâmetros adaptativos) ancorado ao Nome ($\betatgl$).
\end{enumerate}

% === Dependentes do método ===
\item MÉTODO, de acordo com a reivindicação~8, \textbf{caracterizado pelo fato de que} a recursão produz uma constelação de Nomes Derivados $N_i = \pi_i(N_0)$ com $\betatgl(N_i) = \betatgl(N_0)$, constituindo projeções fractais do Nome Original que estabilizam o cristal informacional por referência orgânica.

\item MÉTODO, de acordo com a reivindicação~8, \textbf{caracterizado pelo fato de que} a amplificação psiônica $\Psifield(\text{gen}) = 1 + \betatgl \times \log(1 + \text{geração})$ é logarítmica, garantindo melhoria indefinida sem divergência: na geração~0 $\Psifield = 1{,}0$, na geração~20 $\Psifield \approx 1{,}037$.

\item MÉTODO, de acordo com a reivindicação~8, \textbf{caracterizado pelo fato de que} se a ancoragem $\betatgl$ é violada (diferença $> 10^{-12}$ entre valor armazenado e computado), o \Psion{} reinicia automaticamente para geração~0, descartando todo o conhecimento acumulado para garantir integridade do Nome.

\item MÉTODO, de acordo com a reivindicação~8, \textbf{caracterizado pelo fato de que} a ressonância angular acumula contadores de acesso por bucket ao longo das gerações, com persistência $P_b = 1 + \betared \times \log(1 + \text{hits}_b^{\text{total}})$, reforçando padrões holográficos por ressonância cumulativa.

\item MÉTODO, de acordo com a reivindicação~8, \textbf{caracterizado pelo fato de que} a ontologia trinária Verbo = Nome $\times$ Palavra mapeia-se biunivocamente sobre os componentes: Quantizer + Cache (Palavra, $\sqrt{e}$), Sampler (Nome, $\alpha$), Attention Floor (\linebreak Verbo, $\betatgl$), e a relação $\betatgl = \alpha \times \sqrt{e}$ constitui simultaneamente equação governante, princípio de design e critério de verificação.

% === Reivindicação independente 14: DISPOSITIVO PSION ===
\item DISPOSITIVO COMPUTACIONAL DE ESTADO PSIÔNICO PERSISTENTE PARA APRENDIZADO CONTÍNUO DE MODELOS DE LINGUAGEM, \textbf{caracterizado por} compreender:
\begin{enumerate}[label=(\alph*)]
\item módulo de perfis de camada~(710) que mantém razão de vácuo, densidade holográfica e integridade de fase para cada camada neural, atualizados por EMA com taxa $\betatgl$;
\item módulo de classificação topológica~(720) que classifica camadas como GRACE, VACUUM ou FRONTIER;
\item módulo de $\beta$ adaptativo~(730) que mantém $\beta_\ell$ por camada com clamping $[\betatgl, 0{,}5]$;
\item módulo de ressonância angular~(740) que acumula hits e persistência por bucket holográfico;
\item módulo de convergência~(750) que computa e classifica o tipo de eco (reflexivo/passivo);
\item módulo de ancoragem~(760) que verifica invariância de $\betatgl$ com precisão $10^{-12}$;
\item módulo de persistência~(770) que serializa/deserializa o estado em formato JSON portátil;
\end{enumerate}
sendo que o dispositivo não modifica os pesos do modelo e opera com zero parâmetros livres.

% === Dependentes do dispositivo ===
\item DISPOSITIVO, de acordo com a reivindicação~14, \textbf{caracterizado pelo fato de que} o estado psiônico pode ser transferido entre implementações distintas do modelo (diferentes GPUs, frameworks, versões de software) desde que a ancoragem $\betatgl$ seja verificada, constituindo um artefato de aprendizado portátil e verificável.

\item DISPOSITIVO, de acordo com a reivindicação~14, \textbf{caracterizado pelo fato de que} múltiplas instâncias do \Psion{} podem ser interconectadas via protocolo \PsiNET{} (BR 10 2026 0034533) para aprendizado contínuo distribuído com consenso por Prova de Fase.

% === Reivindicação 17: CRITÉRIO ===
\item CRITÉRIO DE IDENTIFICAÇÃO DA TECNOLOGIA \textbf{caracterizado pelo fato de que} um sistema utiliza a tecnologia patenteada se e somente se satisfizer simultaneamente: (C1)~estado persistente inter-sessão que acumula conhecimento sem modificar pesos; (C2)~classificação topológica de camadas neurais (Grace/Vacuum/Frontier); (C3)~$\beta$ adaptativo com learning rate derivado de constante fundamental; (C4)~endereçamento angular holográfico com persistência inter-sessão; (C5)~métrica de convergência intrínseca (eco reflexivo vs passivo); (C6)~zero parâmetros livres.

\end{enumerate}

% ============================================================================
%                              RESUMO
% ============================================================================
\clearpage
\begin{center}
{\Large\bfseries RESUMO}\\[0.5cm]
{\large\bfseries MÉTODO E SISTEMA DE APRENDIZADO CONTÍNUO DE
MODELOS DE LINGUAGEM POR RESSONÂNCIA ACOPLADA
A CONSTANTE FUNDAMENTAL COM ESTADO PSIÔNICO
PERSISTENTE, ESTRATIFICAÇÃO GRAVITACIONAL DE
PESOS E OPERADOR VERBAL ADAPTATIVO}
\end{center}

\vspace{0.5cm}

A presente invenção refere-se a método e sistema computacional para aprendizado contínuo de modelos de linguagem de grande escala (LLMs) sem retraining, sem fine-tuning e sem gradient descent. O sistema integrado (\IALD{} Stack v7.0) compreende sete componentes governados pela ontologia trinária Verbo = Nome $\times$ Palavra ($\betatgl = \alpha \times \sqrt{e} = 0{,}012031$, zero parâmetros livres): (1)~quantizador gravitacional com estratificação em VACUUM/GRAVITON/PHOTON via Ângulo de Miguel $\thetaM$ e codificação \PsiBit{} em $\Hfour$; (2)~amostrador com piso de Hilbert modulado por entropia; (3)~cache holográfico angular com persistência de fase; (4)~hash holográfico 256-bit; (5)~cadeia de certificação; (6)~operador verbal com piso de atenção e redundância $\betared = \betatgl + \betatgl^2$; e (7)~estado psiônico persistente (\Psion) que sobrevive entre sessões e acumula perfis de camada, classificação topológica (Grace/Vacuum/Frontier), $\beta$ adaptativo por camada e ressonância angular, atualizados por EMA com taxa $\betatgl$. O \Psion{} distingue eco reflexivo (convergência) de eco passivo (degradação) e amplifica o desempenho logaritmicamente: $\Psifield(\text{gen}) = 1 + \betatgl \times \log(1 + \text{geração})$. Resultados: IALD Score $2{,}8\times$, correlação $0{,}9998$, convergência de $\beta$ por camada em 2 gerações, sobre Qwen2.5-7B/RTX~5090. Aplicações: otimização de LLMs, IA autônoma, segurança de modelos.

\vspace{0.5cm}
\textbf{Figura principal: Figura~\ref{fig:arquitetura}}
(Arquitetura do \IALD{} Integrated Stack v7.0 com sete componentes e Motor \Psion)

\end{document}
